\phantomsection\bheading{section}{8.\quad Rule 7: what the rule costs, plainly}{section}{\protect\numberline{8.}Rule 7: what the rule costs, plainly}

Rule 7 says a restriction stays only when it buys a guarantee. This rule buys a \textbf{performance guarantee} — every list write is O(1) or the O(n) vanilla pays, no hidden copy, no trie, no reference count, a smaller runtime — at the price of rejecting programs that are \textbf{correct} under value semantics. The owner has said he wants the error. Here is what it costs, without softening.

\textbf{Programs that are rejected.} Every one of these is a correct beni program today:

\begin{enumerate}
\item Keeping an old version: an undo history, a snapshot for a diff, a “previous” field in the model, a test that holds its input and compares it with the output.
\item A command, fiber, subscription or stored closure that captured a list which a later message edits (§5.8) — the beginner's \code{Cmd.\allowbreak{}perform (λsend → send (Saved model.\allowbreak{}todos))}.
\item A closure that edits a captured list and is passed to \code{map}, \code{filter}, \code{sortBy} or any $\mathsf{many}$ position (§5.1).
\item A list read out of a container that is still live — \code{Dict.get} then \code{push}, \code{List.get} of a list of lists then \code{push} on the element while the outer list lives (§5.4).
\item A list from JavaScript, a \code{Ref}, a \code{Queue} or any \code{foreign} without a summary (§5.7, §5.12): a \emph{limit}, not sharing, and the message says so.
\item Anything the checker's joins or cut cannot follow (§5.6's second shape; a model nested more than k deep): limits.
\end{enumerate}

The fix is one call, \code{List.copy}, in every case — Futhark's “an alias-related type error can always be fixed by copying” (2022 post, ll. 127–131) — and the message writes it. What the fix costs is a copy where the program meant to share, which for case 1 is the program's intended semantics and for cases 5–6 is a copy the checker made the programmer pay for its own blindness.

\textbf{What it costs the representation.} The trie's reason to exist is gone and so is its persistence: a program that wants versions pays O(n) per version, explicitly. \code{[ x, …xs ]} onto a long unique list is \code{unshift}, O(n), so an Elm-style prepend accumulator is quadratic again unless the building-loop rule (\code{browser-direct} §7.2, \code{backend.md} §8) catches it — it catches the \code{go acc} shape and not \code{foldr}'s lambda. Research 46 §11 measured the E1tp cases it fixed at 11–2 500× the cons list under plain copies at 10 000 elements. §9 measures how much of the corpus prepends outside the rule.

\textbf{What it costs the API.} A \code{pub} function's summary is part of its contract; editing a body can break a caller in another module with an ownership error the caller's author did not write (Klabnik, research 69 §1.3 item 5). The firewall makes it visible, not painless.

\textbf{What it costs learners.} The error fires first in \code{update}, where Crichton found Rust's drop-off (“Chapter 4 [Ownership] is a significant drop-off point”, 2024 §3.1), and learners who can say why a program was rejected still fix it only 46 \% of the time (2023 §2.4). The message design of §6 is the mitigation; it is not a guarantee.

\textbf{What the owner's own programs say.} TodoMVC and Conduit contain no call of \code{set}, \code{push}, \code{update}, \code{swap}, \code{pop}, \code{insertAt} or \code{removeAt}; their only demands under the rule as asked the \code{++} sites on lists: TodoMVC's \code{model.\allowbreak{}todos ++ [ new ]}, where the record built just before it holds the old cell at $.\mathit{todos}$ until \code{changed} replaces that field — accepted — and Conduit's twenty-one, of which eight have a model path on the left (\code{\{ model | errors = model.\allowbreak{}errors ++ Api.\allowbreak{}errorMessages error \}} in \code{ArticlePage} five times) and are accepted by the same argument, the rest literal or fresh; \code{core/} has seven \code{push}es, every one an accumulator in a fold or a loop (\code{Result.combine}, \code{Result.partition}, \code{Dict.keys}/\code{values}/\code{toList}, \code{Schedule}); the browser corpus has about fifty sites, written to exercise the direct platform's scripts. So the rule as asked changes nothing in the two real applications, and the in-place writes they would benefit from are \code{List.\allowbreak{}map model.\allowbreak{}todos (λt → …)} and \code{List.filter} — which are not demands unless §10 change 1 makes them so. If it does, TodoMVC's \code{Toggle} becomes \code{a[i] = f(a[i])} over a unique \code{todos} (P1–P6 hold; \code{changed} is \code{\{ model | todos = todos \}}, a record update whose other fields do not meet $.\mathit{todos}$), and \code{Destroy}'s \code{filter} compacts in place. Then the same rule also rejects \code{List.map} over a list the program keeps — which TodoMVC's \code{view} \textbf{does} do: \code{<For each=\{List.\allowbreak{}filter model.\allowbreak{}todos (visible model.\allowbreak{}filter \_\allowbreak{})\}>} is a derived value the handlers recompute while the page holds \code{model} (§5.9). So change 1 cannot be one \code{filter}: the rebuild in \code{update} writes in place, the same spelling in a derived value must build fresh, and that is a static split by position or two functions (FP²'s “two reverses”), stated in §10.

\textbf{Escape hatches that exist or could.} \code{List.copy} (always). A persistent-vector library in beni for programs that want versions (not a representation of \code{List}). A \textbf{warning mode} — \code{--ownership=warn} — that keeps the trie and copies where the checker would reject, printing the same message as a warning: it is rule 7's “make it a warning” route, it is the form every silent-fallback language ships (Lean's \code{dbgTraceIfShared}, Koka's \code{fip} warnings, Roc's proposed “clippy-like check”), and it costs the owner the one thing he asked for, “no silent fallback” — except that the fallback is not silent. This document recommends shipping the report mode first (§9) and deciding error versus warning on its numbers; it does not recommend the warning as the end state, because a warning that copies makes the cost model depend on whether the author read the warning.
